%*************************************
% ภาคผนวก ง.
%*************************************
\appendixChapterTitle{ง}{กระบวนการทำงาน NETCONF/YANG ในระบบ}
\vspace*{-0.8in}

\section{ภาพรวมโปรโตคอล NETCONF}
\hspace*{1.3em}
NETCONF (Network Configuration Protocol) เป็นโปรโตคอลมาตรฐานตาม RFC 6241 ที่ใช้ในการจัดการคอนฟิกอุปกรณ์เครือข่ายผ่านช่องทาง SSH ที่ปลอดภัย โดยใช้ XML เป็นรูปแบบข้อมูลสำหรับการสื่อสารระหว่าง Client (ระบบเว็บแอปพลิเคชัน) กับ Server (อุปกรณ์เครือข่าย) ในโครงงานนี้ระบบใช้ NETCONF เป็นกลไกหลักในการเชื่อมต่อและสั่งงานอุปกรณ์ Cisco Nexus (NX-OS) และ IOS-XE

\hspace*{1.3em}
โปรโตคอล NETCONF มีสถาปัตยกรรมแบบ 4 ชั้น (Layer) ดังนี้

\begin{mycustomenum2}
    \item \textbf{Transport Layer:} ใช้ SSH (Secure Shell) เป็นช่องทางการสื่อสาร โดยเชื่อมต่อผ่าน SSH Subsystem ชื่อ \textbf{netconf} ที่พอร์ต 830 (ค่าเริ่มต้นมาตรฐาน)
    \item \textbf{Messages Layer:} ใช้ XML-based RPC (Remote Procedure Call) ในการส่งคำสั่งและรับผลลัพธ์ โดยข้อความทุกข้อความสิ้นสุดด้วย Delimiter \textbf{]]>]]>} เพื่อระบุจุดสิ้นสุดของข้อความ
    \item \textbf{Operations Layer:} กำหนดชุดคำสั่ง (Operations) ที่สามารถดำเนินการได้ เช่น \textbf{get-config}, \textbf{edit-config}, \textbf{lock}, \textbf{unlock}, \textbf{commit}, \textbf{validate} และ \textbf{close-session}
    \item \textbf{Content Layer:} ข้อมูลคอนฟิกที่จัดการผ่าน NETCONF จะถูกกำหนดโครงสร้างด้วย YANG Data Model
\end{mycustomenum2}

% ──────────────────────────────────────────────────────
% รูปภาพ: NETCONF Protocol Stack / Architecture
% ──────────────────────────────────────────────────────
\begin{figure}[H]
\centering
\adjustbox{frame, width=1.0\textwidth}{
    \includegraphics[width=1.0\textwidth]{Image/netconf-architecture.png}
}
\caption{\fontSixTeen{สถาปัตยกรรมของโปรโตคอล NETCONF แบบ 4 ชั้น}}
\label{figD:NetconfArchitecture}
\end{figure}

\section{YANG Data Model}
\hspace*{1.3em}
YANG (Yet Another Next Generation) เป็นภาษามาตรฐานตาม RFC 7950 สำหรับกำหนดโครงสร้างข้อมูล (Data Model) ของอุปกรณ์เครือข่าย YANG กำหนดว่าข้อมูลคอนฟิกที่จะส่งผ่าน NETCONF ต้องมีรูปแบบอย่างไร ประกอบด้วยฟิลด์อะไรบ้าง แต่ละฟิลด์มีชนิดข้อมูลอะไร และมีข้อจำกัดอย่างไร

\hspace*{1.3em}
ในโครงงานนี้รองรับ YANG Namespace ของอุปกรณ์ Cisco 2 ประเภทหลัก ดังนี้

\begin{table}[H]
\centering
\caption{\fontSixTeen{YANG Namespace และ Root Element สำหรับอุปกรณ์ Cisco ที่รองรับในระบบ}}
\label{tabD:YangNamespaces}
\begin{tabularx}{\textwidth}{|l|X|l|}
\hline
\textbf{ประเภทอุปกรณ์} & \textbf{YANG Namespace} & \textbf{Root Element} \\
\hline
Cisco NX-OS (Nexus) & \textbf{http://cisco.com/ns/yang/cisco-nx-os-device} & \textbf{System} \\
\hline
Cisco IOS-XE & \textbf{http://cisco.com/ns/yang/Cisco-IOS-XE-native} & \textbf{native} \\
\hline
\end{tabularx}
\end{table}

\hspace*{1.3em}
โครงสร้าง YANG Model ประกอบด้วยองค์ประกอบหลักดังนี้

\begin{mycustomenum2}
    \item \textbf{module:} หน่วยหลักของ YANG Model กำหนดชื่อโมดูลและ Namespace
    \item \textbf{namespace:} URI ที่ระบุตัวตนของโมดูล ใช้ใน XML เพื่อระบุว่าข้อมูลเป็นของโมดูลใด
    \item \textbf{prefix:} ชื่อย่อสำหรับอ้างอิง Namespace ในเอกสาร XML
    \item \textbf{container:} โหนดที่ใช้จัดกลุ่มข้อมูลย่อย (เปรียบเสมือนโฟลเดอร์)
    \item \textbf{list:} โหนดที่มีข้อมูลหลายรายการ โดยแต่ละรายการใช้ \textbf{key} ในการระบุตัวตน
    \item \textbf{leaf:} โหนดปลายสุดที่เก็บค่าข้อมูลจริง เช่น ชื่ออินเทอร์เฟซ, IP Address
    \item \textbf{leaf-list:} โหนดปลายสุดที่เก็บค่าได้หลายค่า (เช่น รายการ VLAN ที่อนุญาต)
\end{mycustomenum2}

% ──────────────────────────────────────────────────────
% รูปภาพ: YANG Model Structure / Tree
% ──────────────────────────────────────────────────────
\begin{figure}[H]
\centering
\adjustbox{frame, width=1.0\textwidth}{
    \includegraphics[width=1.0\textwidth]{Image/yang-model-structure.png}
}
\caption{\fontSixTeen{โครงสร้าง YANG Data Model แสดงความสัมพันธ์ระหว่าง Module, Container, List และ Leaf}}
\label{figD:YangModelStructure}
\end{figure}

\hspace*{1.3em}
ในระบบนี้ผู้ใช้สามารถอัปโหลด YANG Model เองได้ โดยระบบจะ Parse ข้อมูลจาก YANG Content อัตโนมัติ ดึงค่า \textbf{module}, \textbf{namespace}, \textbf{prefix} และ \textbf{revision} มาจัดเก็บในฐานข้อมูล นอกจากนี้แต่ละ YANG Model สามารถมี XML Templates เพื่อใช้เป็น Template สำหรับ LLM ในการสร้าง Configuration ได้

\section{กระบวนการเชื่อมต่อ NETCONF Session}
\hspace*{1.3em}
การเชื่อมต่อ NETCONF Session ในระบบมีขั้นตอนดังต่อไปนี้

\subsection{ขั้นตอนที่ 1: SSH Connection}
\hspace*{1.3em}
ระบบเริ่มต้นด้วยการสร้าง SSH Connection ไปยังอุปกรณ์เครือข่ายที่พอร์ต NETCONF (ค่าเริ่มต้นคือ 830) โดยใช้ SSH2 Library ใน Node.js พร้อมกำหนด Algorithms ที่เข้ากันได้กับอุปกรณ์ Cisco ดังนี้

\begin{mycustomenum2}
    \item \textbf{Key Exchange:} \textbf{diffie-hellman-group14-sha256}, \textbf{ecdh-sha2-nistp256}
    \item \textbf{Cipher:} \textbf{aes256-ctr}, \textbf{aes128-ctr}
    \item \textbf{Connection Timeout:} 30 วินาที
\end{mycustomenum2}

\subsection{ขั้นตอนที่ 2: NETCONF Subsystem}
\hspace*{1.3em}
เมื่อ SSH Connection สำเร็จ ระบบจะเปิด NETCONF Subsystem ผ่านคำสั่ง \textbf{conn.subsys('netconf')} เพื่อเข้าสู่โหมดการสื่อสาร NETCONF โดยเฉพาะ ซึ่งจะได้ Stream สำหรับรับ-ส่งข้อความ XML

\subsection{ขั้นตอนที่ 3: Hello Exchange}
\hspace*{1.3em}
ขั้นตอนแรกของโปรโตคอล NETCONF คือการแลกเปลี่ยน Hello Message ระหว่าง Client กับ Server เพื่อตกลง Capabilities ที่ทั้งสองฝ่ายรองรับ

\hspace*{1.3em}
\textbf{Server Hello:} อุปกรณ์เครือข่ายจะส่ง Hello Message มาก่อน ระบบจะ Parse เพื่อดึง Capabilities ที่อุปกรณ์รองรับ เช่น \textbf{base:1.0}, \textbf{writable-running}, \textbf{candidate} เป็นต้น

\hspace*{1.3em}
\textbf{Client Hello:} ระบบตอบกลับด้วย Hello Message ที่ประกาศ Capabilities ที่ Client รองรับ ตัวอย่าง Hello Message ที่ระบบส่งไปมีดังนี้

\begin{verbatim}
<?xml version="1.0" encoding="UTF-8"?>
<hello xmlns="urn:ietf:params:xml:ns:netconf:base:1.0">
  <capabilities>
    <capability>
      urn:ietf:params:netconf:base:1.0
    </capability>
    <capability>
      urn:ietf:params:netconf:capability:writable-running:1.0
    </capability>
    <capability>
      urn:ietf:params:netconf:capability:candidate:1.0
    </capability>
    <capability>
      urn:ietf:params:netconf:capability:confirmed-commit:1.0
    </capability>
    <capability>
      urn:ietf:params:netconf:capability:validate:1.0
    </capability>
  </capabilities>
</hello>]]>]]>
\end{verbatim}

\hspace*{1.3em}
เมื่อทั้งสองฝ่ายแลกเปลี่ยน Hello Message สำเร็จ NETCONF Session จะถูกสร้างขึ้นและพร้อมสำหรับการดำเนินการ RPC ระบบจะจัดเก็บ Session ไว้ใน \textbf{connections Map} พร้อมข้อมูล Capabilities, Stream และ Timestamp เพื่อใช้ในการดำเนินการต่อไป

% ──────────────────────────────────────────────────────
% รูปภาพ: NETCONF Hello Exchange Sequence Diagram
% ──────────────────────────────────────────────────────
\begin{figure}[H]
\centering
\adjustbox{frame, width=1.0\textwidth}{
    \includegraphics[width=1.0\textwidth]{Image/netconf-hello-exchange.png}
}
\caption{\fontSixTeen{Sequence Diagram แสดงกระบวนการ Hello Exchange ระหว่าง Client และ Server}}
\label{figD:HelloExchange}
\end{figure}

\section{NETCONF RPC Operations ที่ใช้ในระบบ}
\hspace*{1.3em}
ระบบใช้ NETCONF RPC Operations หลายชนิดในการจัดการคอนฟิกอุปกรณ์เครือข่าย ทุก RPC Message จะถูกห่อด้วย \textbf{<rpc>} Envelope พร้อมระบุ \textbf{message-id} ที่ไม่ซ้ำกัน และสิ้นสุดด้วย Delimiter \textbf{]]>]]>} โดยรูปแบบทั่วไปของ RPC Message มีดังนี้

\begin{verbatim}
<?xml version="1.0" encoding="UTF-8"?>
<rpc message-id="msg-1709648000000"
     xmlns="urn:ietf:params:xml:ns:netconf:base:1.0">
  <!-- Operation content here -->
</rpc>]]>]]>
\end{verbatim}

\hspace*{1.3em}
รายละเอียดของ RPC Operations แต่ละชนิดมีดังนี้

\subsection{get-config (อ่านค่าคอนฟิก)}
\hspace*{1.3em}
ใช้สำหรับอ่านค่า Running Configuration จากอุปกรณ์ สามารถเลือกอ่านทั้งหมดหรือเฉพาะส่วนที่ต้องการด้วย Subtree Filter ตัวอย่างการอ่านค่าคอนฟิกทั้งหมดจาก Running Datastore มีดังนี้

\begin{verbatim}
<rpc message-id="msg-001"
     xmlns="urn:ietf:params:xml:ns:netconf:base:1.0">
  <get-config>
    <source>
      <running/>
    </source>
  </get-config>
</rpc>]]>]]>
\end{verbatim}

\hspace*{1.3em}
ตัวอย่างการอ่านเฉพาะส่วน Interface ของอุปกรณ์ NX-OS ด้วย Subtree Filter มีดังนี้

\begin{verbatim}
<rpc message-id="msg-002"
     xmlns="urn:ietf:params:xml:ns:netconf:base:1.0">
  <get-config>
    <source>
      <running/>
    </source>
    <filter type="subtree">
      <System xmlns="http://cisco.com/ns/yang/
               cisco-nx-os-device">
        <intf-items/>
      </System>
    </filter>
  </get-config>
</rpc>]]>]]>
\end{verbatim}

\subsection{edit-config (แก้ไขคอนฟิก)}
\hspace*{1.3em}
ใช้สำหรับแก้ไขค่า Configuration ของอุปกรณ์ สามารถเลือก Target Datastore เป็น \textbf{running} (แก้ไขทันที) หรือ \textbf{candidate} (แก้ไขเป็นร่างก่อน) และกำหนด Default Operation เป็น \textbf{merge} (รวมกับค่าเดิม) หรือ \textbf{replace} (แทนที่ค่าเดิม)

\hspace*{1.3em}
ตัวอย่างการตั้งค่า Hostname ของอุปกรณ์ NX-OS ผ่าน edit-config มีดังนี้

\begin{verbatim}
<rpc message-id="msg-003"
     xmlns="urn:ietf:params:xml:ns:netconf:base:1.0">
  <edit-config>
    <target>
      <running/>
    </target>
    <default-operation>merge</default-operation>
    <config>
      <System xmlns="http://cisco.com/ns/yang/
               cisco-nx-os-device">
        <name>Switch-Core-01</name>
      </System>
    </config>
  </edit-config>
</rpc>]]>]]>
\end{verbatim}

\hspace*{1.3em}
ตัวอย่างการตั้งค่า Interface ของอุปกรณ์ IOS-XE ผ่าน edit-config มีดังนี้

\begin{verbatim}
<rpc message-id="msg-004"
     xmlns="urn:ietf:params:xml:ns:netconf:base:1.0">
  <edit-config>
    <target>
      <running/>
    </target>
    <default-operation>merge</default-operation>
    <config>
      <native xmlns="http://cisco.com/ns/yang/
               Cisco-IOS-XE-native">
        <interface>
          <GigabitEthernet>
            <name>1</name>
            <ip>
              <address>
                <primary>
                  <address>192.168.1.1</address>
                  <mask>255.255.255.0</mask>
                </primary>
              </address>
            </ip>
          </GigabitEthernet>
        </interface>
      </native>
    </config>
  </edit-config>
</rpc>]]>]]>
\end{verbatim}

% ──────────────────────────────────────────────────────
% รูปภาพ: ตัวอย่างหน้าจอระบบแสดง NETCONF XML Configuration
% ──────────────────────────────────────────────────────
\begin{figure}[H]
\centering
\adjustbox{frame, width=1.0\textwidth}{
    \includegraphics[width=1.0\textwidth]{Image/netconf-edit-config-example.png}
}
\caption{\fontSixTeen{ตัวอย่างหน้าจอระบบแสดงการสร้าง NETCONF XML Configuration ด้วย LLM}}
\label{figD:EditConfigExample}
\end{figure}

\subsection{lock / unlock (ล็อกและปลดล็อก Datastore)}
\hspace*{1.3em}
ใช้สำหรับล็อก Datastore เพื่อป้องกันไม่ให้ Session อื่นแก้ไขค่าในขณะที่กำลังทำงาน และปลดล็อกเมื่อดำเนินการเสร็จสิ้น

\begin{verbatim}
<!-- Lock -->
<rpc message-id="msg-005"
     xmlns="urn:ietf:params:xml:ns:netconf:base:1.0">
  <lock>
    <target>
      <candidate/>
    </target>
  </lock>
</rpc>]]>]]>

<!-- Unlock -->
<rpc message-id="msg-006"
     xmlns="urn:ietf:params:xml:ns:netconf:base:1.0">
  <unlock>
    <target>
      <candidate/>
    </target>
  </unlock>
</rpc>]]>]]>
\end{verbatim}

\subsection{commit (ยืนยันการเปลี่ยนแปลง)}
\hspace*{1.3em}
ใช้สำหรับยืนยัน (Commit) การเปลี่ยนแปลงจาก Candidate Datastore ไปยัง Running Datastore ทำให้คอนฟิกมีผลบังคับใช้จริง

\begin{verbatim}
<rpc message-id="msg-007"
     xmlns="urn:ietf:params:xml:ns:netconf:base:1.0">
  <commit/>
</rpc>]]>]]>
\end{verbatim}

\subsection{validate (ตรวจสอบความถูกต้อง)}
\hspace*{1.3em}
ใช้ตรวจสอบว่าคอนฟิกที่เตรียมไว้ใน Candidate Datastore มีความถูกต้องตาม YANG Schema หรือไม่ ก่อนที่จะทำการ Commit

\begin{verbatim}
<rpc message-id="msg-008"
     xmlns="urn:ietf:params:xml:ns:netconf:base:1.0">
  <validate>
    <source>
      <candidate/>
    </source>
  </validate>
</rpc>]]>]]>
\end{verbatim}

\subsection{discard-changes (ยกเลิกการเปลี่ยนแปลง)}
\hspace*{1.3em}
ใช้สำหรับยกเลิกการเปลี่ยนแปลงที่อยู่ใน Candidate Datastore กลับไปเป็นค่าเดิมของ Running Datastore ใช้ในกรณีที่ตรวจพบข้อผิดพลาดก่อนที่จะ Commit

\begin{verbatim}
<rpc message-id="msg-009"
     xmlns="urn:ietf:params:xml:ns:netconf:base:1.0">
  <discard-changes/>
</rpc>]]>]]>
\end{verbatim}

\subsection{close-session (ปิด NETCONF Session)}
\hspace*{1.3em}
ใช้สำหรับปิด NETCONF Session อย่างสมบูรณ์เมื่อดำเนินการเสร็จสิ้น

\begin{verbatim}
<rpc message-id="msg-010"
     xmlns="urn:ietf:params:xml:ns:netconf:base:1.0">
  <close-session/>
</rpc>]]>]]>
\end{verbatim}

\hspace*{1.3em}
ตารางที่ \ref{tabD:RpcOperations} สรุป RPC Operations ทั้งหมดที่ใช้ในระบบพร้อมค่า Timeout ที่กำหนด

\begin{table}[H]
\centering
\caption{\fontSixTeen{สรุป NETCONF RPC Operations ที่ใช้ในระบบ}}
\label{tabD:RpcOperations}
\begin{tabularx}{\textwidth}{|l|l|l|X|}
\hline
\textbf{Operation} & \textbf{RPC} & \textbf{Timeout} & \textbf{คำอธิบาย} \\
\hline
เชื่อมต่อ & Hello Exchange & 30 วินาที & สร้าง SSH + NETCONF Session \\
\hline
อ่านคอนฟิก & \textbf{get-config} & 180 วินาที & อ่าน Running Configuration \\
\hline
อ่านข้อมูล & \textbf{get} & 120 วินาที & อ่าน Operational Data \\
\hline
แก้ไขคอนฟิก & \textbf{edit-config} & 90 วินาที & แก้ไข Running/Candidate \\
\hline
ยืนยัน & \textbf{commit} & 120 วินาที & ยืนยัน Candidate $\rightarrow$ Running \\
\hline
ตรวจสอบ & \textbf{validate} & 60 วินาที & ตรวจสอบ Candidate \\
\hline
ยกเลิก & \textbf{discard-changes} & 15 วินาที & ยกเลิก Candidate \\
\hline
ล็อก/ปลดล็อก & \textbf{lock/unlock} & 30 วินาที & ล็อก Datastore \\
\hline
ปิด Session & \textbf{close-session} & 10 วินาที & ปิด NETCONF Session \\
\hline
\end{tabularx}
\end{table}

\section{กระบวนการ Deploy Configuration ผ่าน NETCONF}
\hspace*{1.3em}
การ Deploy Configuration ผ่าน NETCONF ในระบบมี 2 รูปแบบ ขึ้นอยู่กับ Capabilities ที่อุปกรณ์รองรับ ดังนี้

\subsection{รูปแบบที่ 1: ใช้ Candidate Datastore (แนะนำ)}
\hspace*{1.3em}
วิธีนี้เป็นวิธีที่ปลอดภัยที่สุดตามมาตรฐาน RFC 6241 เนื่องจากการเปลี่ยนแปลงจะถูกเขียนไปยัง Candidate Datastore ก่อน และจะมีผลบังคับใช้ก็ต่อเมื่อทำการ Commit สำเร็จเท่านั้น หากเกิดข้อผิดพลาดสามารถ Discard Changes เพื่อย้อนกลับได้ ขั้นตอนมีดังนี้

\begin{mycustomenum2}
    \item \textbf{Lock Candidate:} ล็อก Candidate Datastore เพื่อป้องกัน Session อื่นแก้ไข
    \item \textbf{Edit-Config:} ส่ง Configuration XML ไปยัง Candidate Datastore ด้วย \textbf{edit-config}
    \item \textbf{Validate (ถ้ารองรับ):} ตรวจสอบความถูกต้องของ Candidate ด้วย \textbf{validate}
    \item \textbf{Commit:} ยืนยันการเปลี่ยนแปลงจาก Candidate ไปยัง Running ด้วย \textbf{commit}
    \item \textbf{Unlock Candidate:} ปลดล็อก Candidate Datastore
\end{mycustomenum2}

\hspace*{1.3em}
หากขั้นตอนใดเกิดข้อผิดพลาด ระบบจะทำ Cleanup อัตโนมัติ ได้แก่ Discard Changes เพื่อยกเลิกการเปลี่ยนแปลงใน Candidate และ Unlock เพื่อปลดล็อก Datastore

% ──────────────────────────────────────────────────────
% รูปภาพ: Candidate Datastore Workflow
% ──────────────────────────────────────────────────────
\begin{figure}[H]
\centering
\adjustbox{frame, width=1.0\textwidth}{
    \includegraphics[width=1.0\textwidth]{Image/netconf-candidate-workflow.png}
}
\caption{\fontSixTeen{Flowchart แสดงกระบวนการ Deploy ผ่าน Candidate Datastore พร้อม Error Handling}}
\label{figD:CandidateWorkflow}
\end{figure}

\subsection{รูปแบบที่ 2: ใช้ Writable-Running (Fallback)}
\hspace*{1.3em}
สำหรับอุปกรณ์ที่ไม่รองรับ Candidate Datastore ระบบจะใช้ \textbf{writable-running} เป็นทางเลือก โดยการเปลี่ยนแปลงจะมีผลทันทีเมื่อ edit-config สำเร็จ ขั้นตอนมีดังนี้

\begin{mycustomenum2}
    \item \textbf{Edit-Config:} ส่ง Configuration XML ไปยัง Running Datastore โดยตรง
    \item การเปลี่ยนแปลงมีผลบังคับใช้ทันที ไม่ต้อง Commit
\end{mycustomenum2}

\hspace*{1.3em}
\textbf{ข้อควรระวัง:} วิธีนี้ไม่มีขั้นตอน Validate ก่อน Commit และไม่สามารถ Rollback ได้หากเกิดข้อผิดพลาด จึงควรใช้ Candidate Datastore เป็นลำดับแรกเมื่อเป็นไปได้

\section{Subtree Filter สำหรับอ่านข้อมูลอุปกรณ์}
\hspace*{1.3em}
ระบบใช้ Subtree Filter ใน \textbf{<get>} RPC เพื่ออ่านข้อมูลเฉพาะส่วนที่ต้องการจากอุปกรณ์ โดยไม่ต้องอ่านข้อมูลทั้งหมด ซึ่งช่วยลดปริมาณข้อมูลที่รับ-ส่งและเพิ่มประสิทธิภาพ ระบบรองรับ Filter หลายประเภทตามชนิดของอุปกรณ์ ดังนี้

\subsection{Subtree Filter สำหรับ NX-OS (Nexus)}
\hspace*{1.3em}
YANG Model ของ Cisco NX-OS ใช้ Root Element \textbf{System} ภายใต้ Namespace \textbf{cisco-nx-os-device} ตัวอย่าง Filter ที่ใช้ในระบบมีดังนี้

\begin{table}[H]
\centering
\caption{\fontSixTeen{Subtree Filter สำหรับอุปกรณ์ Cisco NX-OS}}
\label{tabD:NxosFilters}
\begin{tabularx}{\textwidth}{|l|l|X|}
\hline
\textbf{ประเภทข้อมูล} & \textbf{YANG Path} & \textbf{คำอธิบาย} \\
\hline
System Info & \textbf{System/name} & ชื่ออุปกรณ์ (Hostname) \\
\hline
Interfaces & \textbf{System/intf-items} & อินเทอร์เฟซทั้งหมด (Physical, SVI, Port-channel) \\
\hline
VLANs & \textbf{System/bd-items} & Bridge Domain / VLAN \\
\hline
OSPF & \textbf{System/ospf-items} & Routing Protocol OSPF \\
\hline
BGP & \textbf{System/bgp-items} & Routing Protocol BGP \\
\hline
Features & \textbf{System/fm-items} & Feature Manager (เปิด/ปิดฟีเจอร์) \\
\hline
\end{tabularx}
\end{table}

\subsection{Subtree Filter สำหรับ IOS-XE}
\hspace*{1.3em}
YANG Model ของ Cisco IOS-XE ใช้ Root Element \textbf{native} ภายใต้ Namespace \textbf{Cisco-IOS-XE-native} ตัวอย่าง Filter ที่ใช้ในระบบมีดังนี้

\begin{table}[H]
\centering
\caption{\fontSixTeen{Subtree Filter สำหรับอุปกรณ์ Cisco IOS-XE}}
\label{tabD:IosXeFilters}
\begin{tabularx}{\textwidth}{|l|l|X|}
\hline
\textbf{ประเภทข้อมูล} & \textbf{YANG Path} & \textbf{คำอธิบาย} \\
\hline
Hostname & \textbf{native/hostname} & ชื่ออุปกรณ์ \\
\hline
Interfaces & \textbf{native/interface} & อินเทอร์เฟซ GigabitEthernet/Loopback \\
\hline
VLANs & \textbf{native/vlan} & VLAN Configuration \\
\hline
OSPF & \textbf{native/router/ospf} & Routing Protocol OSPF \\
\hline
BGP & \textbf{native/router/bgp} & Routing Protocol BGP \\
\hline
\end{tabularx}
\end{table}

% ──────────────────────────────────────────────────────
% รูปภาพ: ตัวอย่างหน้าจอระบบแสดงข้อมูล NETCONF ที่อ่านได้
% ──────────────────────────────────────────────────────
\begin{figure}[H]
\centering
\adjustbox{frame, width=1.0\textwidth}{
    \includegraphics[width=1.0\textwidth]{Image/netconf-device-details.png}
}
\caption{\fontSixTeen{ตัวอย่างหน้าจอระบบแสดงข้อมูลอุปกรณ์ที่อ่านผ่าน NETCONF}}
\label{figD:DeviceDetails}
\end{figure}

\section{การ Deploy ผ่าน Agent Relay}
\hspace*{1.3em}
เนื่องจากระบบ Backend ทำงานเป็น Vercel Serverless Functions ซึ่งไม่สามารถเชื่อมต่อ SSH/NETCONF ไปยังอุปกรณ์ในเครือข่ายภายในได้โดยตรง ระบบจึงใช้ Desktop Agent เป็นตัวกลาง (Relay) ในการส่งคำสั่ง NETCONF ไปยังอุปกรณ์ กระบวนการมีดังนี้

\begin{mycustomenum2}
    \item \textbf{Backend สร้าง Command:} เมื่อผู้ใช้สั่ง Deploy ระบบจะสร้าง AgentCommand ใน MongoDB พร้อมข้อมูล NETCONF XML Configuration
    \item \textbf{Agent Poll Command:} Desktop Agent ทำ HTTP Polling ทุก 2 วินาทีเพื่อดึง Pending Commands
    \item \textbf{Agent ดำเนินการ:} Agent เชื่อมต่อ SSH/NETCONF ไปยังอุปกรณ์จริงในเครือข่ายภายใน ดำเนินการตาม Workflow (Lock $\rightarrow$ Edit-Config $\rightarrow$ Validate $\rightarrow$ Commit $\rightarrow$ Unlock)
    \item \textbf{Agent รายงานผล:} Agent ส่งผลลัพธ์กลับมายัง Backend ผ่าน HTTP POST
    \item \textbf{Backend ส่งผลให้ผู้ใช้:} Backend อัปเดตสถานะ Command และแจ้งผลลัพธ์ให้ผู้ใช้ทราบ
\end{mycustomenum2}

% ──────────────────────────────────────────────────────
% รูปภาพ: Agent Relay NETCONF Workflow
% ──────────────────────────────────────────────────────
\begin{figure}[H]
\centering
\adjustbox{frame, width=1.0\textwidth}{
    \includegraphics[width=1.0\textwidth]{Image/netconf-agent-relay-workflow.png}
}
\caption{\fontSixTeen{Sequence Diagram แสดงกระบวนการ Deploy NETCONF Configuration ผ่าน Agent Relay}}
\label{figD:AgentRelayWorkflow}
\end{figure}

\section{การจัดการ YANG Model ในระบบ}
\hspace*{1.3em}
ระบบมีระบบจัดการ YANG Model ที่ช่วยให้ผู้ใช้สามารถอัปโหลด จัดเก็บ และใช้ YANG Model ในการสร้าง Configuration ผ่าน LLM ได้ โดยมีองค์ประกอบหลักดังนี้

\subsection{การอัปโหลดและ Parse YANG Model}
\hspace*{1.3em}
เมื่อผู้ใช้อัปโหลด YANG Model ระบบจะ Parse เนื้อหาอัตโนมัติเพื่อดึงข้อมูลสำคัญ ได้แก่

\begin{mycustomenum2}
    \item \textbf{module name:} ชื่อโมดูล YANG เช่น \textbf{cisco-nx-os-device}
    \item \textbf{namespace:} URI ของโมดูล เช่น \textbf{http://cisco.com/ns/yang/cisco-nx-os-device}
    \item \textbf{prefix:} ชื่อย่อสำหรับอ้างอิง เช่น \textbf{nxos}
    \item \textbf{revision:} เวอร์ชันของ YANG Model เช่น \textbf{2024-01-01}
\end{mycustomenum2}

\subsection{XML Templates}
\hspace*{1.3em}
แต่ละ YANG Model สามารถมี XML Templates หลายรายการ ซึ่งเป็นตัวอย่าง XML Configuration ที่ถูกต้องตาม YANG Schema ระบบจะส่ง Templates เหล่านี้ให้ LLM เป็นตัวอย่างประกอบเพื่อช่วยให้ LLM สร้าง Configuration ที่มี Syntax และ Namespace ถูกต้อง

\subsection{การใช้ YANG Model สำหรับ LLM Generation}
\hspace*{1.3em}
เมื่อผู้ใช้สร้าง Configuration ด้วย LLM ระบบจะ

\begin{mycustomenum2}
    \item ดึง YANG Models ที่ตรงกับประเภทอุปกรณ์ (NX-OS หรือ IOS-XE) จากฐานข้อมูล
    \item จัดรูปแบบข้อมูล YANG Model เป็นข้อความบริบท (Context) ประกอบด้วย Namespace, Config Paths และ XML Templates
    \item ส่งข้อมูลบริบทพร้อม Prompt ของผู้ใช้ไปให้ LLM เพื่อสร้าง Configuration ที่ถูกต้องตาม YANG Schema
\end{mycustomenum2}

% ──────────────────────────────────────────────────────
% รูปภาพ: หน้าจอจัดการ YANG Model
% ──────────────────────────────────────────────────────
\begin{figure}[H]
\centering
\adjustbox{frame, width=1.0\textwidth}{
    \includegraphics[width=1.0\textwidth]{Image/yang-model-management.png}
}
\caption{\fontSixTeen{ตัวอย่างหน้าจอระบบจัดการ YANG Model}}
\label{figD:YangModelManagement}
\end{figure}

\section{การจัดการ Error ใน NETCONF}
\hspace*{1.3em}
เมื่ออุปกรณ์พบข้อผิดพลาดในการดำเนินการ NETCONF จะตอบกลับด้วย \textbf{<rpc-error>} ระบบจะ Parse ข้อมูลใน Error Message เพื่อแสดงรายละเอียดให้ผู้ใช้ทราบ โดย Error Message ประกอบด้วย

\begin{mycustomenum2}
    \item \textbf{error-tag:} ประเภทของข้อผิดพลาด เช่น \textbf{invalid-value}, \textbf{data-missing}, \textbf{operation-failed}
    \item \textbf{error-message:} คำอธิบายข้อผิดพลาดเป็นข้อความภาษาอังกฤษ
    \item \textbf{error-path:} ตำแหน่งใน YANG Model ที่เกิดข้อผิดพลาด เช่น \textbf{/System/intf-items}
\end{mycustomenum2}

\hspace*{1.3em}
ตัวอย่าง RPC Error Response มีดังนี้

\begin{verbatim}
<rpc-reply xmlns="urn:ietf:params:xml:ns:netconf:base:1.0"
           message-id="msg-003">
  <rpc-error>
    <error-type>application</error-type>
    <error-tag>invalid-value</error-tag>
    <error-severity>error</error-severity>
    <error-path>/System/intf-items</error-path>
    <error-message xml:lang="en">
      Invalid interface name
    </error-message>
  </rpc-error>
</rpc-reply>]]>]]>
\end{verbatim}

\hspace*{1.3em}
ในกระบวนการ Deploy ผ่าน Candidate Datastore หากเกิดข้อผิดพลาดตรงขั้นตอนใดก็ตาม ระบบจะทำ Cleanup อัตโนมัติ ประกอบด้วย Discard Changes เพื่อยกเลิกการเปลี่ยนแปลงใน Candidate, Unlock เพื่อปลดล็อก Datastore และแจ้ง Error กลับให้ผู้ใช้ทราบพร้อมรายละเอียด เพื่อให้สามารถแก้ไข Configuration และลองใหม่ได้

% ──────────────────────────────────────────────────────
% รูปภาพ: ตัวอย่างหน้าจอแสดง Error จาก NETCONF
% ──────────────────────────────────────────────────────
\begin{figure}[H]
\centering
\adjustbox{frame, width=1.0\textwidth}{
    \includegraphics[width=1.0\textwidth]{Image/netconf-error-handling.png}
}
\caption{\fontSixTeen{ตัวอย่างหน้าจอแสดงข้อผิดพลาดจากการ Deploy ผ่าน NETCONF}}
\label{figD:ErrorHandling}
\end{figure}
